\section{Counting favorable events in probability}

Combinatorics becomes part of probability when it is used to count sample spaces and events.

If \(\Omega\) is finite and all elementary outcomes are equally likely, then
\[
P(A)=\frac{|A|}{|\Omega|}.
\]
This formula is valid only after the sample space has been correctly chosen and the elementary outcomes are equally likely.

\begin{examplebox}
A fair die is rolled twice. If order is recorded, then
\[
\Omega=\{1,2,3,4,5,6\}^2,\qquad |\Omega|=36.
\]
The event that the sum is \(7\) is
\[
A=\{(1,6),(2,5),(3,4),(4,3),(5,2),(6,1)\}.
\]
Thus
\[
P(A)=\frac{6}{36}=\frac{1}{6}.
\]
\end{examplebox}

\begin{examplebox}
Five cards are drawn from a standard deck, and the final hand is recorded. If every five-card hand is equally likely, then
\[
|\Omega|=\binom{52}{5}.
\]
The number of hands with exactly two hearts is
\[
\binom{13}{2}\binom{39}{3}.
\]
Therefore
\[
P(\text{exactly two hearts})=
\frac{\binom{13}{2}\binom{39}{3}}{\binom{52}{5}}.
\]
\end{examplebox}

The important point is that the denominator and numerator must be counted inside the same model. It is wrong to count ordered outcomes in the denominator and unordered outcomes in the numerator, or conversely.

\section{Changing the recording rule changes the model}

The same experiment may be counted in different ways because different questions require different sample spaces.

\begin{examplebox}
Three balls are drawn without replacement from an urn containing labelled balls. Compare the following recording rules.

\begin{center}
\begin{tabular}{lll}
\toprule
Recorded information & Outcome type & Typical count \\
\midrule
labels in order & sequence without repetition & \(P(n,3)\) \\
labels only, order ignored & subset & \(\binom{n}{3}\) \\
colors in order & sequence of colors & depends on color counts \\
color counts only & count vector & depends on available colors \\
number of red balls only & value of a random variable & \(0,1,2,3\) possible \\
\bottomrule
\end{tabular}
\end{center}
\end{examplebox}

The physical action is the same: three balls are drawn. The mathematical model is not the same, because the recorded outcome is different.
